\documentclass[11pt,a4paper]{article}
\usepackage[T1]{fontenc}
\usepackage{lmodern,amsmath,amssymb,amsthm,mathtools,microtype}
\usepackage[margin=25mm]{geometry}
\usepackage[colorlinks=true,linkcolor=blue,citecolor=blue,urlcolor=blue]{hyperref}
\setlength{\emergencystretch}{2em}
\newtheorem{theorem}{Theorem}[section]
\newtheorem{lemma}[theorem]{Lemma}
\newtheorem{corollary}[theorem]{Corollary}
\newcommand{\Q}{\mathbb Q}\newcommand{\Z}{\mathbb Z}
\DeclareMathOperator{\rad}{rad}
\title{Simple primitive denominator primes under weak \texorpdfstring{$abc$}{abc}}
\author{}\date{17 November 2026}
\begin{document}\maketitle

\begin{abstract}
For a fixed congruent-number curve over the rationals and a non-torsion
point in its doubled subgroup, weak rational $abc$ with exponent $16/15$
implies that every sufficiently large denominator term has a primitive
prime of valuation one, outside any prescribed finite set.
Consequently, inclusion of the odd denominator support of one sufficiently
large term in another denominator forces divisibility of the indices.
The proof uses a quartic identity and canonical height, without a
Diophantine approximation estimate for the archimedean denominator.
\end{abstract}

\section{Denominator heights}

Fix an integer $d\ge1$ and the elliptic curve
\[
                       E_d:y^2=x^3-d^2x.
\]
Let $P\in2E_d(\Q)$ have infinite order. For $n\ge1$, write
\begin{equation}\label{eq:coordinates}
 x([n]P)=\frac{A_n^2}{B_n^2},\qquad
 y([n]P)=\frac{A_n C_n}{B_n^3},\qquad
 A_n,B_n>0,\quad\gcd(A_n,B_n)=1.
\end{equation}
These expressions exist because
\[
 x([2]R)=\left(\frac{x(R)^2+d^2}{2y(R)}\right)^2
\]
for a non-torsion point $R$. The equation of the curve then gives
\begin{equation}\label{eq:quartic}
                         C_n^2=A_n^4-d^2B_n^4.
\end{equation}
Initially $C_n$ is rational; its square is integral, so $C_n\in\Z$.
It is nonzero, since $[n]P$ is not torsion. In particular
$A_n^4>d^2B_n^4$, and
\[
                  h_n:=h(A_n:B_n)=\log A_n.
\]
Normalize the canonical height by
$\widehat h(R)=\tfrac12h(x(R))+O_{E_d}(1)$. Its standard bounded
comparison with the naive height and its quadraticity give
\begin{equation}\label{eq:canonical}
              h_n=n^2\widehat h(P)+O_{E_d}(1),\qquad
                         \widehat h(P)>0.
\end{equation}
These height facts are as in \cite[Chapter VIII, Section 9]{AEC}.

For a finite set $S$ of primes, define the simple denominator support
\[
                 L_{n,S}=\prod_{\substack{p\notin S\\v_p(B_n)=1}}p.
\]
Assume weak rational $abc$ with exponent $\kappa>1$: there is
$K_\kappa\ge1$ such that
\[
 c\le K_\kappa\rad(abc)^\kappa
\]
for coprime positive integers $a+b=c$.

\begin{lemma}\label{lem:radical}
Uniformly in $n$,
\[
 \log\rad(B_n)\ge\left(\frac4\kappa-3\right)h_n-O_{d,\kappa,K_\kappa}(1),
\]
and
\begin{equation}\label{eq:simple}
 \log L_{n,S}\ge\left(\frac8\kappa-7\right)h_n
                                      -O_{d,S,\kappa,K_\kappa}(1).
\end{equation}
\end{lemma}
\begin{proof}
Put $g=\gcd(A_n^4,d^2B_n^4)$. Since $A_n,B_n$ are coprime,
$g=\gcd(A_n^4,d^2)\le d^2$. Dividing
$C_n^2+d^2B_n^4=A_n^4$ by $g$ gives a coprime positive $abc$ triple.
Its prime support is contained in that of $dA_nB_nC_n$. Therefore
\[
 \frac{A_n^4}{g}\le K_\kappa\rad(dA_nB_nC_n)^\kappa
                  \le K_\kappa\bigl(d A_n^3\rad(B_n)\bigr)^\kappa,
\]
where $|C_n|<A_n^2$ was used. This proves the first estimate.
For any positive integer $B$, prime-by-prime comparison gives
\[
       \prod_{v_p(B)=1}p\ \ge\ \frac{\rad(B)^2}{B}.
\]
Use $\log B_n\le h_n$ and omit the fixed finite set $S$, at a
bounded logarithmic cost, to obtain \eqref{eq:simple}.
\end{proof}

\section{Primitive primes and index divisibility}

A prime $p$ is primitive for $B_n$ if $p\mid B_n$ and
$p\nmid B_m$ for every $1\le m<n$. At a prime of good reduction,
\begin{equation}\label{eq:order}
 p\mid B_n\quad\Longleftrightarrow\quad [n]\widetilde P=O
                    \quad\Longleftrightarrow\quad r_p\mid n,
\end{equation}
where $r_p$ is the order of $\widetilde P\in E_d(\mathbf F_p)$.
This follows directly from the reduction homomorphism: negative
valuation of the affine abscissa means reduction to the identity.

The usual primitive-denominator and odd-primitivity results are
discussed in \cite[Lemmas 9 and 13]{Silverman} and
\cite[Proposition 3.10 and the following remarks]{CZ}.
The following statement specifies a weak exponent and valuation one.

\begin{theorem}\label{thm:primitive}
Assume weak rational $abc$ with exponent $16/15$.
For every fixed finite set $S$ of primes, there exists $n_0$ such
that, for every $n\ge n_0$, the integer $B_n$ has a primitive prime
$p\notin S$ with $v_p(B_n)=1$.
\end{theorem}
\begin{proof}
Enlarge $S$ by the primes dividing $2d$. Every prime outside $S$
is then a prime of good reduction. If $p\mid B_n$ is not primitive,
\eqref{eq:order} gives a proper divisor $r_p<n$ of $n$. Choose a
prime $\ell\mid n/r_p$. Then $r_p\mid n/\ell$, so
$p\mid B_{n/\ell}$. Thus the product of all nonprimitive prime
divisors of $B_n$ outside $S$ divides
\[
                              \prod_{\ell\mid n\atop\ell\ {
m prime}}B_{n/\ell}.
\]
If there is no simple primitive prime outside $S$, this gives
\begin{align*}
 \log L_{n,S}
 &\le\sum_{\ell\mid n\atop\ell\ {
m prime}}\log B_{n/\ell}\\
 &\le n^2\widehat h(P)\sum_{\ell\mid n\atop\ell\ {
m prime}}\ell^{-2}
                                                     +O_{E_d}(\omega(n)).
\end{align*}
An elementary bound, sufficient here, is
\begin{align*}
 \sum_{\ell\ {
m prime}}\ell^{-2}
 &\le\frac14+\frac19+\frac1{25}
                    +\sum_{j=3}^{\infty}(2j+1)^{-2}\\
 &\le\frac14+\frac19+\frac1{25}+\frac1{49}
                         +\int_3^{\infty}(2x+1)^{-2}\,dx\\
 &=\frac{21739}{44100}<\frac12.
\end{align*}
On the other hand $8/(16/15)-7=1/2$, so
Lemma~\ref{lem:radical} and \eqref{eq:canonical} imply
\[
 \frac12 n^2\widehat h(P)-O(1)\le\log L_{n,S}.
\]
Consequently
\[
                    \frac{311}{44100}n^2\widehat h(P)
                                       \le O(1)+O(\omega(n)),
\]
which bounds $n$, since $\omega(n)\le\log_2 n$.
This contradiction proves the theorem for every sufficiently large $n$.
\end{proof}

\begin{corollary}\label{cor:divides}
Under the same hypothesis, fix a finite set $S$ containing the bad
reduction primes. For every pair of nonzero integers $h,j$ with
$|h|\ge n_0$, the implication
\[
 \bigl(\forall p\notin S:\ v_p(B_{|h|})\text{ odd }
                                      \Longrightarrow p\mid B_{|j|}\bigr)
                          \quad\Longrightarrow\quad h\mid j
\]
holds. Conversely, $h\mid j$ implies the displayed support inclusion.
\end{corollary}
\begin{proof}
Take the simple primitive prime of $B_{|h|}$ from
Theorem~\ref{thm:primitive}. Its reduction order is $|h|$.
The assumed support inclusion puts it in $B_{|j|}$, so
\eqref{eq:order} gives $|h|\mid |j|$. The reverse implication
follows from \eqref{eq:order} at every prime outside $S$.
\end{proof}

The theorem concerns curves over $\Q$ and rational $abc$.
It makes no assertion over a quadratic number field from rational
$abc$ alone. The support-inclusion predicate in Corollary~\ref{cor:divides}
is not given an existential encoding here. Nor is an equality between
a rational input and the index ratio $j/h$ established. These are
additional requirements for an application to a defining polynomial.

\begin{thebibliography}{9}\raggedright
\bibitem{AEC}J.~H.~Silverman, \emph{The Arithmetic of Elliptic Curves},
Graduate Texts in Mathematics 106, second edition, Springer, 2009,
Chapter VIII, Section 9.
\bibitem{Silverman}J.~H.~Silverman,
\emph{Wieferich's criterion and the $abc$-conjecture},
Journal of Number Theory 30 (1988), 226--237,
\url{https://doi.org/10.1016/0022-314X(88)90019-4}.
\bibitem{CZ}G.~Cornelissen and K.~Zahidi,
\emph{Elliptic divisibility sequences and undecidable problems about
rational points}, Journal f\"ur die reine und angewandte Mathematik
613 (2007), 1--33, arXiv:math/0412473v3,
\url{https://arxiv.org/html/math/0412473v3}.
\end{thebibliography}
\end{document}
